\documentclass[tikz,border=8pt]{standalone}
\usepackage{tikz}
\usepackage{amsmath}
\usepackage{amssymb}
\usepackage{pgfplots}
\pgfplotsset{compat=1.18}
\usepackage{ctex}
\usetikzlibrary{calc,positioning,arrows.meta,trees,matrix}

% 排列回溯决策树（LC 46 Permutations）
% nums=[1,2,3]，完整决策树，6 个叶子

\begin{document}
\begin{tikzpicture}[
  every node/.style={},
  edge from parent/.style={draw,->,>=Stealth,thick},
  level 1/.style={sibling distance=5.6cm, level distance=2.0cm},
  level 2/.style={sibling distance=2.0cm, level distance=2.0cm},
  level 3/.style={sibling distance=1.8cm, level distance=2.0cm},
]

%% ── 标题 ────────────────────────────────────────────────────
\node[font=\large\bfseries] at (0, 2.5) {回溯决策树：全排列（LC 46 Permutations）};
\node[font=\small,gray] at (0, 1.9) {nums = [1, 2, 3]，蓝色 = 递归进入，红色 = 回溯撤销};

%% ── 根节点 ───────────────────────────────────────────────────
\node[draw,rectangle,rounded corners=2pt,minimum width=1.2cm,minimum height=0.6cm,
      fill=blue!10,font=\bfseries] (root) at (0,0) {[\ ]}
  child {
    node[draw,rectangle,rounded corners=2pt,minimum width=1.2cm,minimum height=0.6cm,
         fill=blue!15,font=\bfseries] (n1) {[1]}
    child {
      node[draw,rectangle,rounded corners=2pt,minimum width=1.4cm,minimum height=0.6cm,
           fill=blue!20,font=\bfseries] (n12) {[1,2]}
      child {
        node[draw,rectangle,rounded corners=2pt,minimum width=1.6cm,minimum height=0.6cm,
             fill=green!25,font=\bfseries] (n123) {[1,2,3]}
      }
    }
    child {
      node[draw,rectangle,rounded corners=2pt,minimum width=1.4cm,minimum height=0.6cm,
           fill=blue!20,font=\bfseries] (n13) {[1,3]}
      child {
        node[draw,rectangle,rounded corners=2pt,minimum width=1.6cm,minimum height=0.6cm,
             fill=green!25,font=\bfseries] (n132) {[1,3,2]}
      }
    }
  }
  child {
    node[draw,rectangle,rounded corners=2pt,minimum width=1.2cm,minimum height=0.6cm,
         fill=blue!15,font=\bfseries] (n2) {[2]}
    child {
      node[draw,rectangle,rounded corners=2pt,minimum width=1.4cm,minimum height=0.6cm,
           fill=blue!20,font=\bfseries] (n21) {[2,1]}
      child {
        node[draw,rectangle,rounded corners=2pt,minimum width=1.6cm,minimum height=0.6cm,
             fill=green!25,font=\bfseries] (n213) {[2,1,3]}
      }
    }
    child {
      node[draw,rectangle,rounded corners=2pt,minimum width=1.4cm,minimum height=0.6cm,
           fill=blue!20,font=\bfseries] (n23) {[2,3]}
      child {
        node[draw,rectangle,rounded corners=2pt,minimum width=1.6cm,minimum height=0.6cm,
             fill=green!25,font=\bfseries] (n231) {[2,3,1]}
      }
    }
  }
  child {
    node[draw,rectangle,rounded corners=2pt,minimum width=1.2cm,minimum height=0.6cm,
         fill=blue!15,font=\bfseries] (n3) {[3]}
    child {
      node[draw,rectangle,rounded corners=2pt,minimum width=1.4cm,minimum height=0.6cm,
           fill=blue!20,font=\bfseries] (n31) {[3,1]}
      child {
        node[draw,rectangle,rounded corners=2pt,minimum width=1.6cm,minimum height=0.6cm,
             fill=green!25,font=\bfseries] (n312) {[3,1,2]}
      }
    }
    child {
      node[draw,rectangle,rounded corners=2pt,minimum width=1.4cm,minimum height=0.6cm,
           fill=blue!20,font=\bfseries] (n32) {[3,2]}
      child {
        node[draw,rectangle,rounded corners=2pt,minimum width=1.6cm,minimum height=0.6cm,
             fill=green!25,font=\bfseries] (n321) {[3,2,1]}
      }
    }
  };

%% ── 层级标注（左侧） ──────────────────────────────────────────
\node[font=\scriptsize,gray,align=right] at (-9.5, 0)   {第 0 层\\（选择数为 3）};
\node[font=\scriptsize,gray,align=right] at (-9.5,-2.0) {第 1 层\\（选择数为 2）};
\node[font=\scriptsize,gray,align=right] at (-9.5,-4.0) {第 2 层\\（选择数为 1）};
\node[font=\scriptsize,gray,align=right] at (-9.5,-6.0) {第 3 层\\（叶子/结果）};

%% ── 边标注（选择 + 回溯） ─────────────────────────────────────
% root → n1
\node[font=\scriptsize,blue!70!black] at (-5.9,-0.85) {选 1};
\node[font=\scriptsize,red!70!black]  at (-5.3,-1.35) {\textit{撤}};
% root → n2
\node[font=\scriptsize,blue!70!black] at (0.1,-0.85) {选 2};
\node[font=\scriptsize,red!70!black]  at (0.7,-1.35) {\textit{撤}};
% root → n3
\node[font=\scriptsize,blue!70!black] at (5.7,-0.85) {选 3};
\node[font=\scriptsize,red!70!black]  at (6.3,-1.35) {\textit{撤}};

% n1 → n12
\node[font=\scriptsize,blue!70!black] at (-7.1,-2.85) {选 2};
% n1 → n13
\node[font=\scriptsize,blue!70!black] at (-4.7,-2.85) {选 3};
% n2 → n21
\node[font=\scriptsize,blue!70!black] at (-1.5,-2.85) {选 1};
% n2 → n23
\node[font=\scriptsize,blue!70!black] at (1.5,-2.85)  {选 3};
% n3 → n31
\node[font=\scriptsize,blue!70!black] at (4.3,-2.85) {选 1};
% n3 → n32
\node[font=\scriptsize,blue!70!black] at (6.8,-2.85) {选 2};

%% ── 回溯弧（红色虚线）连叶子回到第1层 ────────────────────────
\draw[->,>=Stealth,red!60,dashed,bend right=30,thick]
  (n123.south east) to node[font=\tiny,red!60,right] {回溯} (n12.east);
\draw[->,>=Stealth,red!60,dashed,bend left=30,thick]
  (n132.south west) to node[font=\tiny,red!60,left]  {回溯} (n13.west);
\draw[->,>=Stealth,red!60,dashed,bend right=30,thick]
  (n213.south east) to node[font=\tiny,red!60,right] {回溯} (n21.east);
\draw[->,>=Stealth,red!60,dashed,bend left=30,thick]
  (n231.south west) to node[font=\tiny,red!60,left]  {回溯} (n23.west);
\draw[->,>=Stealth,red!60,dashed,bend right=30,thick]
  (n312.south east) to node[font=\tiny,red!60,right] {回溯} (n31.east);
\draw[->,>=Stealth,red!60,dashed,bend left=30,thick]
  (n321.south west) to node[font=\tiny,red!60,left]  {回溯} (n32.west);

%% ── 叶子说明 ────────────────────────────────────────────────
\node[draw=gray!50,fill=green!8,rounded corners=3pt,font=\small,inner sep=6pt,align=center]
  at (0,-8.0)
  {\textbf{6 种排列结果（叶子）}\\[4pt]
   [1,2,3]\quad [1,3,2]\quad [2,1,3]\quad [2,3,1]\quad [3,1,2]\quad [3,2,1]};

%% ── 算法说明 ────────────────────────────────────────────────
\node[draw=gray!50,fill=white,rounded corners=3pt,font=\small,inner sep=6pt,align=left]
  at (0,-9.5)
  {\textbf{回溯模板：}
   \texttt{for x in 剩余选择: path.append(x); backtrack(); path.pop()}\\
   时间复杂度 $O(n \cdot n!)$，空间 $O(n)$（递归栈）};

\end{tikzpicture}
\end{document}
